\section{Empirical context}\label{sec:empirical}

The Six Birds vocabulary used in this paper was first tested in a simulation study of a stochastic substrate. That study, the PICA study~\cite{Tsiokos2026ToCreate}, organizes closure mechanisms as a $6\times 6$ table of interaction cells between primitives. A configuration, or \emph{program}, switches on a subset of cells. This section summarizes what the study's shipped measurements do and do not show about the primitive-role reading used here. \emph{No theorem in this paper depends on these data.} We report them for context and because they bound the claims one may attach to the vocabulary. All numbers below are recomputed from the study's saved run summaries, and no new simulations were run.

\paragraph{Measurement.}
The outcome is $\mathrm{frob}$, the Frobenius distance between the coarse-grained transition kernel and the rank-one stationary kernel. It is averaged over ten seeds per configuration and system size~$n$. For a configuration $P$ we use the absolute shift $|\Delta(P)|=|\overline{\mathrm{frob}}(P)-\overline{\mathrm{frob}}(\text{baseline})|$. The \emph{gain} of $P$ over a comparator $Q$ is $|\Delta(P)|-|\Delta(Q)|$. Two comparators are used. The \emph{nearest} comparator is the largest nonempty proper subprogram of $P$ in the study's atlas, with a fixed tie-break. The \emph{best} comparator is the proper subprogram with the largest absolute shift, among those that appear in the study's atlas, including the empty program. A configuration that genuinely needs all of its interactions should beat both.

\begin{figure}[t]
\centering
% Generated by scripts/make_paper_figures.py from
% results/ticket-p6-1/interaction_delta_table.csv -- do not edit by hand.
\begin{tikzpicture}[font=\footnotesize]
\begin{scope}[xshift=0.00cm]
\draw[black!12, line width=0.4pt] (0.000,-0.25) -- (0.000,4.570);
\node[below, black!70, font=\scriptsize] at (0.000,-0.25) {-0.8};
\draw[black!12, line width=0.4pt] (0.900,-0.25) -- (0.900,4.570);
\node[below, black!70, font=\scriptsize] at (0.900,-0.25) {-0.6};
\draw[black!12, line width=0.4pt] (1.800,-0.25) -- (1.800,4.570);
\node[below, black!70, font=\scriptsize] at (1.800,-0.25) {-0.4};
\draw[black!12, line width=0.4pt] (2.700,-0.25) -- (2.700,4.570);
\node[below, black!70, font=\scriptsize] at (2.700,-0.25) {-0.2};
\draw[black!55, line width=0.6pt] (3.600,-0.25) -- (3.600,4.570);
\node[below, black!70, font=\scriptsize] at (3.600,-0.25) {0};
\draw[black!12, line width=0.4pt] (4.500,-0.25) -- (4.500,4.570);
\node[below, black!70, font=\scriptsize] at (4.500,-0.25) {+0.2};
\draw[black!12, line width=0.4pt] (5.400,-0.25) -- (5.400,4.570);
\node[below, black!70, font=\scriptsize] at (5.400,-0.25) {+0.4};
\node[font=\footnotesize\bfseries] at (2.700,4.920) {$n=64$};
\draw[black!30, line width=1pt] (3.513,4.320) -- (0.838,4.320);
\node[rectangle, fill=seriesB, draw=white, line width=0.6pt, inner sep=1.9pt] at (0.838,4.320) {};
\node[circle, fill=seriesA, draw=white, line width=0.6pt, inner sep=1.9pt] at (3.513,4.320) {};
\node[anchor=east, font=\scriptsize\ttfamily] at (-0.1,4.320) {A13\_A14};
\draw[black!30, line width=1pt] (4.078,3.960) -- (1.403,3.960);
\node[rectangle, fill=seriesB, draw=white, line width=0.6pt, inner sep=1.9pt] at (1.403,3.960) {};
\node[circle, fill=seriesA, draw=white, line width=0.6pt, inner sep=1.9pt] at (4.078,3.960) {};
\node[anchor=east, font=\scriptsize\ttfamily] at (-0.1,3.960) {A14\_A19};
\draw[black!30, line width=1pt] (3.977,3.600) -- (0.925,3.600);
\node[rectangle, fill=seriesB, draw=white, line width=0.6pt, inner sep=1.9pt] at (0.925,3.600) {};
\node[circle, fill=seriesA, draw=white, line width=0.6pt, inner sep=1.9pt] at (3.977,3.600) {};
\node[anchor=east, font=\scriptsize\ttfamily] at (-0.1,3.600) {A14\_only};
\draw[black!30, line width=1pt] (4.692,3.240) -- (1.640,3.240);
\node[rectangle, fill=seriesB, draw=white, line width=0.6pt, inner sep=1.9pt] at (1.640,3.240) {};
\node[circle, fill=seriesA, draw=white, line width=0.6pt, inner sep=1.9pt] at (4.692,3.240) {};
\node[anchor=east, font=\scriptsize\ttfamily] at (-0.1,3.240) {A16\_only};
\draw[black!30, line width=1pt] (4.044,2.880) -- (0.991,2.880);
\node[rectangle, fill=seriesB, draw=white, line width=0.6pt, inner sep=1.9pt] at (0.991,2.880) {};
\node[circle, fill=seriesA, draw=white, line width=0.6pt, inner sep=1.9pt] at (4.044,2.880) {};
\node[anchor=east, font=\scriptsize\ttfamily] at (-0.1,2.880) {A17\_only};
\draw[black!30, line width=1pt] (4.210,2.520) -- (1.285,2.520);
\node[rectangle, fill=seriesB, draw=white, line width=0.6pt, inner sep=1.9pt] at (1.285,2.520) {};
\node[circle, fill=seriesA, draw=white, line width=0.6pt, inner sep=1.9pt] at (4.210,2.520) {};
\node[anchor=east, font=\scriptsize\ttfamily] at (-0.1,2.520) {A18\_A19};
\draw[black!30, line width=1pt] (3.600,2.160) -- (0.548,2.160);
\node[rectangle, fill=seriesB, draw=white, line width=0.6pt, inner sep=1.9pt] at (0.548,2.160) {};
\node[circle, fill=seriesA, draw=white, line width=0.6pt, inner sep=1.9pt] at (3.600,2.160) {};
\node[anchor=east, font=\scriptsize\ttfamily] at (-0.1,2.160) {A20\_only};
\draw[black!30, line width=1pt] (3.754,1.800) -- (0.701,1.800);
\node[rectangle, fill=seriesB, draw=white, line width=0.6pt, inner sep=1.9pt] at (0.701,1.800) {};
\node[circle, fill=seriesA, draw=white, line width=0.6pt, inner sep=1.9pt] at (3.754,1.800) {};
\node[anchor=east, font=\scriptsize\ttfamily] at (-0.1,1.800) {P1\_row};
\draw[black!30, line width=1pt] (2.877,1.440) -- (0.562,1.440);
\node[rectangle, fill=seriesB, draw=white, line width=0.6pt, inner sep=1.9pt] at (0.562,1.440) {};
\node[circle, fill=seriesA, draw=white, line width=0.6pt, inner sep=1.9pt] at (2.877,1.440) {};
\node[anchor=east, font=\scriptsize\ttfamily] at (-0.1,1.440) {P3\_row};
\draw[black!30, line width=1pt] (3.600,1.080) -- (0.548,1.080);
\node[rectangle, fill=seriesB, draw=white, line width=0.6pt, inner sep=1.9pt] at (0.548,1.080) {};
\node[circle, fill=seriesA, draw=white, line width=0.6pt, inner sep=1.9pt] at (3.600,1.080) {};
\node[anchor=east, font=\scriptsize\ttfamily] at (-0.1,1.080) {P5\_row};
\draw[black!30, line width=1pt] (3.977,0.720) -- (0.924,0.720);
\node[rectangle, fill=seriesB, draw=white, line width=0.6pt, inner sep=1.9pt] at (0.924,0.720) {};
\node[circle, fill=seriesA, draw=white, line width=0.6pt, inner sep=1.9pt] at (3.977,0.720) {};
\node[anchor=east, font=\scriptsize\ttfamily] at (-0.1,0.720) {P6\_row};
\draw[black!30, line width=1pt] (4.612,0.360) -- (1.703,0.360);
\node[rectangle, fill=seriesB, draw=white, line width=0.6pt, inner sep=1.9pt] at (1.703,0.360) {};
\node[circle, fill=seriesA, draw=white, line width=0.6pt, inner sep=1.9pt] at (4.612,0.360) {};
\node[anchor=east, font=\scriptsize\ttfamily] at (-0.1,0.360) {full\_action};
\draw[black!30, line width=1pt] (2.791,0.000) -- (0.894,0.000);
\node[rectangle, fill=seriesB, draw=white, line width=0.6pt, inner sep=1.9pt] at (0.894,0.000) {};
\node[circle, fill=seriesA, draw=white, line width=0.6pt, inner sep=1.9pt] at (2.791,0.000) {};
\node[anchor=east, font=\scriptsize\ttfamily] at (-0.1,0.000) {full\_all};
\node[black!70, font=\scriptsize] at (2.700,-0.85) {gain in $|\Delta\,\mathrm{frob}|$ over comparator};
\end{scope}
\begin{scope}[xshift=6.80cm]
\draw[black!12, line width=0.4pt] (0.000,-0.25) -- (0.000,4.570);
\node[below, black!70, font=\scriptsize] at (0.000,-0.25) {-0.8};
\draw[black!12, line width=0.4pt] (0.900,-0.25) -- (0.900,4.570);
\node[below, black!70, font=\scriptsize] at (0.900,-0.25) {-0.6};
\draw[black!12, line width=0.4pt] (1.800,-0.25) -- (1.800,4.570);
\node[below, black!70, font=\scriptsize] at (1.800,-0.25) {-0.4};
\draw[black!12, line width=0.4pt] (2.700,-0.25) -- (2.700,4.570);
\node[below, black!70, font=\scriptsize] at (2.700,-0.25) {-0.2};
\draw[black!55, line width=0.6pt] (3.600,-0.25) -- (3.600,4.570);
\node[below, black!70, font=\scriptsize] at (3.600,-0.25) {0};
\draw[black!12, line width=0.4pt] (4.500,-0.25) -- (4.500,4.570);
\node[below, black!70, font=\scriptsize] at (4.500,-0.25) {+0.2};
\draw[black!12, line width=0.4pt] (5.400,-0.25) -- (5.400,4.570);
\node[below, black!70, font=\scriptsize] at (5.400,-0.25) {+0.4};
\node[font=\footnotesize\bfseries] at (2.700,4.920) {$n=128$};
\draw[black!30, line width=1pt] (3.009,4.320) -- (1.030,4.320);
\node[rectangle, fill=seriesB, draw=white, line width=0.6pt, inner sep=1.9pt] at (1.030,4.320) {};
\node[circle, fill=seriesA, draw=white, line width=0.6pt, inner sep=1.9pt] at (3.009,4.320) {};
\draw[black!30, line width=1pt] (2.636,3.960) -- (0.657,3.960);
\node[rectangle, fill=seriesB, draw=white, line width=0.6pt, inner sep=1.9pt] at (0.657,3.960) {};
\node[circle, fill=seriesA, draw=white, line width=0.6pt, inner sep=1.9pt] at (2.636,3.960) {};
\draw[black!30, line width=1pt] (5.106,3.600) -- (1.621,3.600);
\node[rectangle, fill=seriesB, draw=white, line width=0.6pt, inner sep=1.9pt] at (1.621,3.600) {};
\node[circle, fill=seriesA, draw=white, line width=0.6pt, inner sep=1.9pt] at (5.106,3.600) {};
\draw[black!30, line width=1pt] (4.494,3.240) -- (1.009,3.240);
\node[rectangle, fill=seriesB, draw=white, line width=0.6pt, inner sep=1.9pt] at (1.009,3.240) {};
\node[circle, fill=seriesA, draw=white, line width=0.6pt, inner sep=1.9pt] at (4.494,3.240) {};
\draw[black!30, line width=1pt] (4.343,2.880) -- (0.857,2.880);
\node[rectangle, fill=seriesB, draw=white, line width=0.6pt, inner sep=1.9pt] at (0.857,2.880) {};
\node[circle, fill=seriesA, draw=white, line width=0.6pt, inner sep=1.9pt] at (4.343,2.880) {};
\draw[black!30, line width=1pt] (3.638,2.520) -- (0.153,2.520);
\node[rectangle, fill=seriesB, draw=white, line width=0.6pt, inner sep=1.9pt] at (0.153,2.520) {};
\node[circle, fill=seriesA, draw=white, line width=0.6pt, inner sep=1.9pt] at (3.638,2.520) {};
\draw[black!30, line width=1pt] (3.600,2.160) -- (0.115,2.160);
\node[rectangle, fill=seriesB, draw=white, line width=0.6pt, inner sep=1.9pt] at (0.115,2.160) {};
\node[circle, fill=seriesA, draw=white, line width=0.6pt, inner sep=1.9pt] at (3.600,2.160) {};
\draw[black!30, line width=1pt] (4.213,1.800) -- (0.727,1.800);
\node[rectangle, fill=seriesB, draw=white, line width=0.6pt, inner sep=1.9pt] at (0.727,1.800) {};
\node[circle, fill=seriesA, draw=white, line width=0.6pt, inner sep=1.9pt] at (4.213,1.800) {};
\draw[black!30, line width=1pt] (3.591,1.440) -- (0.144,1.440);
\node[rectangle, fill=seriesB, draw=white, line width=0.6pt, inner sep=1.9pt] at (0.144,1.440) {};
\node[circle, fill=seriesA, draw=white, line width=0.6pt, inner sep=1.9pt] at (3.591,1.440) {};
\draw[black!30, line width=1pt] (3.600,1.080) -- (0.115,1.080);
\node[rectangle, fill=seriesB, draw=white, line width=0.6pt, inner sep=1.9pt] at (0.115,1.080) {};
\node[circle, fill=seriesA, draw=white, line width=0.6pt, inner sep=1.9pt] at (3.600,1.080) {};
\draw[black!30, line width=1pt] (3.806,0.720) -- (0.321,0.720);
\node[rectangle, fill=seriesB, draw=white, line width=0.6pt, inner sep=1.9pt] at (0.321,0.720) {};
\node[circle, fill=seriesA, draw=white, line width=0.6pt, inner sep=1.9pt] at (3.806,0.720) {};
\draw[black!30, line width=1pt] (3.829,0.360) -- (0.661,0.360);
\node[rectangle, fill=seriesB, draw=white, line width=0.6pt, inner sep=1.9pt] at (0.661,0.360) {};
\node[circle, fill=seriesA, draw=white, line width=0.6pt, inner sep=1.9pt] at (3.829,0.360) {};
\draw[black!30, line width=1pt] (3.646,0.000) -- (0.707,0.000);
\node[rectangle, fill=seriesB, draw=white, line width=0.6pt, inner sep=1.9pt] at (0.707,0.000) {};
\node[circle, fill=seriesA, draw=white, line width=0.6pt, inner sep=1.9pt] at (3.646,0.000) {};
\node[black!70, font=\scriptsize] at (2.700,-0.85) {gain in $|\Delta\,\mathrm{frob}|$ over comparator};
\end{scope}
\node[circle, fill=seriesA, inner sep=1.9pt] at (0.2,5.470) {};
\node[anchor=west] at (0.35,5.470) {vs.\ nearest nonempty proper subprogram};
\node[rectangle, fill=seriesB, inner sep=1.9pt] at (6.6,5.470) {};
\node[anchor=west] at (6.75,5.470) {vs.\ best measured proper subprogram};
\end{tikzpicture}

\caption{Gains in absolute $\mathrm{frob}$ shift for the 13 configurations at the two primary system sizes (26 comparisons). Each row joins the gain over the nearest nonempty proper subprogram (blue circles) to the gain over the best measured proper subprogram (orange squares). Points right of zero favor the full configuration. Sixteen nearest-comparator gains are positive, but every best-comparator gain is negative, and in all 26 cases the best comparator is the empty program. Means over ten seeds per cell; descriptive, with no uncertainty estimate.}
\label{fig:comparators}
\end{figure}

\paragraph{What the comparisons show.}
Figure~\ref{fig:comparators} plots both gains for all 26 comparisons at $n\in\{64,128\}$. Against the nearest comparator, 16 gains are positive, 4 are zero and 6 are negative. The largest is $0.335$, for \texttt{A14\_only} at $n=128$. Against the best comparator, all 26 gains are negative, ranging from $-0.42$ to $-0.77$. In every case the best comparator is the empty program, which already shifts $\mathrm{frob}$ further from the baseline than any of the configurations does. A positive nearest-comparator gain is therefore a local comparison with one neighbor. It is not evidence that the full set of interactions is needed. Appendix~\ref{app:empirical-details} lists the values.

\paragraph{A readiness check fixed in advance.}
Before the confirmatory runs, this project froze decision rules for promoting one configuration, \texttt{A14\_only}, to a theorem target. The rules required a \emph{signed} increase $\overline{\mathrm{frob}}(P)-\overline{\mathrm{frob}}(\text{baseline})$ of at least $0.15$ at both $n=64$ and $n=128$, a margin of at least $0.05$ in mean absolute shift over the two neighboring configurations \texttt{A13\_A14} and \texttt{A14\_A19}, complete seed coverage at the anchor resolution $k=4$, and no new robustness failures. Applying these rules to the existing logs gives the verdict \emph{not ready}. All $120$ primary seed-level metrics are present, but $18$ seeds, spread over six configuration--size cells, lack a measured $k=4$ anchor. The signed increase at $n=64$ is $0.084$, below the $0.15$ threshold, and the robustness condition is not established. Only the comparator margin is met: the mean absolute shift of \texttt{A14\_only} over the two primary sizes is $0.209$, against $0.155$ and $0.134$ for its neighbors.

\paragraph{Consequences for this paper.}
The data do not show that richer interaction structure outperforms simpler subprograms, and they do not single out a primitive interaction as necessary. This is consistent with the way the primitives are used here, as interpretive labels attached to proved statements and not as hypotheses that carry proof. The labels make no claim beyond the mathematics. In particular, the empirical record does not license treating protocol mismatch (P3) as a driver of theory growth, and P3 plays no role in any theorem.
